\section{Symmetry}

\noindent\textbf{Exercise~\ref{ex:symmetry-1}.}\label{sol:symmetry-1}

By definition, we have
\bse
    (\phi_*X)\la f \ra = X \la f\circ \phi \ra,
\ese
for $X\in T\cM$, giving
\bse
    \begin{split}
        \phi^*(df) : X & := df : \phi_*(X) \\
        & = \phi_* X \la f \ra \\
        & = X \la f\circ \phi \ra \\
        & = d(f\circ \phi) : X \\
        & =: d(\phi^*f) :X,
    \end{split}
\ese
which holds for arbitrary $X$ and therefore proves the result.

\bigskip
\noindent\textbf{Exercise~\ref{ex:symmetry-2}.}\label{sol:symmetry-2}

If we considered a general vector and gradient, we would have
\bse
    df : \phi_*X = \phi_*X\la f \ra = X\la f \circ \phi \ra = X^i \bigg(\frac{\p(f\circ\phi)}{\p x^i}\bigg).
\ese
Now we have
\bse
    \begin{split}
        \bigg(\frac{\p (f\circ\phi)}{\p x^i}\bigg)_p & := \p_i\big(f\circ\phi\circ x^{-1}\big)\big|_{x(p)} \\
        & = \p_i\big(f\circ y^{-1}\circ y\circ \phi\circ x^{-1}\big)\big|_{x(p)} \\
        & = \p_j\big(f\circ y^{-1}\big)\big|_{(y\circ\phi\circ x^{-1}\circ x)(p)} \cdot \p_i\big(y^i\circ \phi\circ x^{-1}\big)\big|_{x(p)} \\
        & =: \bigg(\frac{\p f}{\p y^j}\bigg)_q \cdot \bigg(\frac{\p (y^j\circ\phi)}{\p x^i}\bigg)_p,
    \end{split}
\ese
where $q:=\phi(p)\in\cN$. Now put in $f=y^a$ and $X=\frac{\p}{\p x^b}$, giving
\bse
    \begin{split}
        \phi_{*\,\, b}^{\,\, a}(p) & = \bigg(\frac{\p y^a}{\p y^j}\bigg)_q \cdot \bigg(\frac{\p (y^j\circ \phi)}{\p x^b}\bigg)_p \\
        & = \del^a_j \cdot \bigg(\frac{\p (y^j\circ \phi)}{\p x^b}\bigg)_p \\
        & = \bigg(\frac{\p (y\circ \phi)^a}{\p x^b}\bigg)_p.
    \end{split}
\ese

\bigskip
\noindent\textbf{Exercise~\ref{ex:symmetry-3}.}\label{sol:symmetry-3}

With the definition of the induced metric in mind, we have
\bse
    \begin{split}
        (\phi^*g)(X,Y) & = g\big(\phi_*X,\phi_*Y\big) \\
        & = g_{ab}\big(\phi_*X\big)^a\big(\phi_*Y\big)^b \\
        & = g_{ab} \big(dy^a:\phi_*X\big) \big(dy^b:\phi_*Y\big).
    \end{split}
\ese
Then if we use $g_{ab}=g\big(\frac{\p}{\p x^a},\frac{\p}{\p x^b}\big)$ and the results of the previous exercise we get
\bse
    \begin{split}
        (\phi^*g)_{ab}(p) & = g_{mn}(q) \cdot \bigg[dy^m:\phi_*\bigg(\frac{\p}{\p x^a}\bigg)\bigg](p) \cdot \bigg[dy^n:\phi_*\bigg(\frac{\p}{\p x^b}\bigg)\bigg](p) \\
        & = \bigg(\frac{\p(y\circ\phi)^m}{\p x^a}\bigg)_p \bigg(\frac{\p(y\circ\phi)^n}{\p x^b}\bigg)_p g_{mn}\big(\phi(p)\big).
    \end{split}
\ese

\bigskip
\noindent\textbf{Exercise~\ref{ex:symmetry-4}.}\label{sol:symmetry-4}

Nothing as in order to talk about shape you need either a covariant derivative or a metric, neither of which we have. I guess you could say it is some closed and compact 2-dimensional shape, but you could not specify which one, i.e. if its a round sphere or an ellipsoid or a potato.

\bigskip
\noindent\textbf{Exercise~\ref{ex:symmetry-5}.}\label{sol:symmetry-5}

Using the result of the previous questions, we need to find
\bse
    \bigg(\frac{\p (y\circ \iota)^m}{\p x^a}\bigg)_p := \p_a\big(y^m\circ \iota \circ x^{-1}\big)\big|_{x(p)}.
\ese
Using the definition given, we have
\bse
    \begin{split}
        \bigg(\frac{\p (y\circ \iota)^1}{\p x^1}\bigg)_p & = a\cos\varphi\cos\vartheta, \qquad \bigg(\frac{\p (y\circ \iota)^1}{\p x^2}\bigg)_p = -a\sin\varphi\sin\vartheta \\
        \bigg(\frac{\p (y\circ \iota)^2}{\p x^1}\bigg)_p & = b\sin\varphi\cos\vartheta, \qquad \bigg(\frac{\p (y\circ \iota)^2}{\p x^2}\bigg)_p  = b\cos\varphi\sin\vartheta \\
        \bigg(\frac{\p (y\circ \iota)^3}{\p x^1}\bigg)_p & = -c\sin\vartheta, \qquad \qquad \bigg(\frac{\p (y\circ \iota)^3}{\p x^2}\bigg)_p = 0.
    \end{split}
\ese
You then just plug in the relevant terms, giving
\bse
    \begin{split}
        g^{\text{ellipsoid}}_{11} & = a^2 \cos^2\varphi\cos^2\vartheta + b^2\sin^2\varphi\cos^2\vartheta + c^2\sin^2\vartheta \\
        g^{\text{ellipsoid}}_{22} & = a^2\sin^2\varphi\sin^2\vartheta + b^2\cos^2\varphi\sin^2\vartheta,
    \end{split}
\ese
and $g^{\text{ellipsoid}}_{12}=0=g^{\text{ellipsoid}}_{21}$.

\bigskip
\noindent\textbf{Exercise~\ref{ex:symmetry-6}.}\label{sol:symmetry-6}

Let $f\in C^{\infty}(\cM)$ be an arbitrary smooth function. We need to consider the action of the Lie bracket expressions on $f$, e.g. $[X_1,X_2]\la f\ra$. We use the clever trick that in this expansion only the terms where a derivative acts on a term in the $X$s will remain. That is, any terms that are second order derivative of $f$ will vanish because it will appear in both $X_1\la X_2\la f\ra \ra$ and $X_2\la X_1\la f\ra \ra$ which the order switched, but partial derivative commute and so these terms cancel.

We then have (dropping the $p$s for notational reasons)
\bse
    \begin{split}
        [X_1,X_2]\la f \ra & := X_1\big\la X_2\la f\ra \big\ra - X_2\big\la X_1\la f\ra \big\ra \\
        & = \bigg[\big(-\cosec^2\vartheta \sin^2\varphi + \cot^2\vartheta\cos^2\varphi\big) \bigg(\frac{\p f}{\p\varphi}\bigg) + \cot\vartheta\cos\varphi\sin\varphi \bigg(\frac{\p f}{\p \vartheta}\bigg)\bigg]  \\
        & \qquad - \bigg[\big(\cosec^2\vartheta\cos^2\varphi - \cot^2\vartheta\sin^2\varphi\big) \bigg(\frac{\p f}{\p \varphi}\bigg) + \cot\vartheta\sin\varphi\cos\varphi \bigg(\frac{\p f}{\p \vartheta}\bigg)\bigg] \\
        & = \big(\cot^2\vartheta-\cosec^2\vartheta\big)\bigg(\frac{\p f}{\p \varphi}\bigg) \\
        & = -\bigg(\frac{\p }{\p \varphi}\bigg)\la f \ra \\
        & = - X_3\la f \ra \\
        \implies [X_2,X_1] & = X_3,
    \end{split}
\ese
where on the last line we have used the antisymmetry of the Lie bracket.

Next we have
\bse
    \begin{split}
        [X_1,X_3]\la f \ra & = 0 - \bigg[-\cos\varphi \bigg(\frac{\p f}{\p \vartheta}\bigg) + \cot\vartheta\sin\varphi\bigg(\frac{\p f}{\p \varphi}\bigg)\bigg] \\
        & = X_2\la f \ra \\
        \implies [X_1,X_3] & = X_2,
    \end{split}
\ese
and
\bse
    \begin{split}
        [X_3,X_2] & = \bigg[ -\sin\varphi\bigg(\frac{\p f}{\p \vartheta}\bigg) -\cot\vartheta\cos\varphi\bigg(\frac{\p f}{\p \varphi}\bigg) \bigg] - 0 \\
        & = X_1\la f\ra \\
        \implies [X_3,X_2] = X_1.
    \end{split}
\ese
So we see that $\{X_1,X_2,X_3\}$ is closed under the Lie bracket, and so forms a Lie subalgebra. The structure constants are
\bse
    {C^3}_{21} = {C^2}_{13} = {C^1}_{32} = 1,
\ese
and all other non-related (i.e. not ${C^3}_{12} = - {C^3}_{21}$, etc.) structure constants vanish.

Note this result tells us that $\{X_1,X_2,X_3\}$ is a 3-dimensional rotation algebra, as defined in the lecture. We therefore expect it to be a symmtry of $S^2$, which we show explictly below for $X_3$.

\bigskip
\noindent\textbf{Exercise~\ref{ex:symmetry-7}.}\label{sol:symmetry-7}

Using
\bse
    v_{\gamma_p,\gamma_p(\lambda)} = (X_3)_{\gamma_p(\lambda)} \qquad \iff \qquad \dot{\gamma}^i_{p(x)}\bigg(\frac{\p}{\p x^i}\bigg)_{\gamma_p(\lambda)} = \bigg(\frac{\p}{\p \varphi}\bigg)_{\gamma_p(\lambda)},
\ese
we have
\bse
    \dot{\gamma}^1_{p(x)}(\lambda) := \big(x^1\circ \gamma_p\big)'(\lambda) = 0, \qand \dot{\gamma}^2_{p(x)}(\lambda) := \big(x^2\circ \gamma_p\big)'(\lambda) = 1,
\ese
from which is follows that
\bse
    \vartheta(\gamma_p) = a, \qand \varphi(\gamma_p) = \lambda + b,
\ese
for constants $a$ and $b$. We see from the question that $a=\vartheta_0$ and $b=\varphi_0$. So we have
\bse
    \gamma_{p(x)}(\lambda) = \big(\vartheta_0, \lambda + \varphi_0\big),
\ese
which satisfies
\bse
    \gamma_p(0) = x^{-1}\big(\gamma_{p(x)}(0)\big) = x^{-1}(\vartheta_0,\varphi_0) = p.
\ese

\bigskip
\noindent\textbf{Exercise~\ref{ex:symmetry-8}.}\label{sol:symmetry-8}

The flow is
\bse
    h^{X_3}_{\lambda} : p \mapsto \gamma_p(\lambda),
\ese
so we have
\bse
    \Big(x^m\circ h^{X_3}_{\lambda}\circ x^{-1}\Big) : (\vartheta,\varphi) \mapsto \gamma^m_p(\lambda),
\ese
and so
\bse
    \begin{split}
        \bigg(\frac{\p (x\circ h^{X_3}_{\lambda})^1}{\p x^1}\bigg)_p & := \p_1\Big( x^1 \circ h^{X_3}_{\lambda}\circ x^{-1}\Big)\Big|_{x(p)} = \vartheta_0 \\
        \bigg(\frac{\p (x\circ h^{X_3}_{\lambda})^1}{\p x^2}\bigg)_p & := \p_2\Big( x^1 \circ h^{X_3}_{\lambda}\circ x^{-1}\Big)\Big|_{x(p)} = 0 \\
        \bigg(\frac{\p (x\circ h^{X_3}_{\lambda})^2}{\p x^1}\bigg)_p & := \p_1\Big( x^2 \circ h^{X_3}_{\lambda}\circ x^{-1}\Big)\Big|_{x(p)} = 0 \\
        \bigg(\frac{\p (x\circ h^{X_3}_{\lambda})^2}{\p x^2}\bigg)_p & := \p_2\Big( x^2 \circ h^{X_3}_{\lambda}\circ x^{-1}\Big)\Big|_{x(p)} = \lambda + \varphi_0.
    \end{split}
\ese
This gives us
\bse
    \begin{split}
        \Big[\Big(h^{X_3}_{\lambda}\Big)^*g^{\text{ellipsoid}}\Big]_{11}(p) & = \vartheta_0^2 g^{\text{ellipsoid}}_{11}\big(\gamma_p(\lambda)\big) \\
        \Big[\Big(h^{X_3}_{\lambda}\Big)^*g^{\text{ellipsoid}}\Big]_{22}(p) & = (\lambda+\varphi_0)^2 g^{\text{ellipsoid}}_{22}\big(\gamma_p(\lambda)\big)
    \end{split}
\ese
and again the other two components vanish. If we then take the coordinate transformation
\bse
    \vartheta \to \frac{1}{\vartheta_0}\vartheta, \qand \varphi \to \frac{1}{\lambda+\varphi_0}\varphi,
\ese
which we can do as $\vartheta_0,\varphi_0 >0$ (as the ranges of $\vartheta$ and $\varphi$ are positive), we then get
\bse
    \Big[\Big(h^{X_3}_{\lambda}\Big)^*g^{\text{ellipsoid}}\Big]_{ab}(p) = g^{\text{ellipsoid}}_{ab}\big(\gamma_p(\lambda)\big),
\ese
or more nicely
\bse
    \Big(h^{X_3}_{\lambda}\Big)^*g^{\text{ellipsoid}} = g^{\text{ellipsoid}}.
\ese
This tells us that $X_3$ is a symmetry of the metric, and so $\cL_{X_3}g^{\text{ellipsoid}} = 0$.
